\chapter{جمع‌بندی و پیشنهادها}
\label{chap:conclusion}

\section{جمع‌بندی}
این گزارش چارچوبی اولیه برای کاربرد شبکه‌های عصبی فیزیک‌آگاه در مدیریت انرژی ریزشبکه ارائه کرد. مسئله مدیریت انرژی مجموعه‌ای از تصمیم‌های به‌هم‌پیوسته درباره تأمین بار، تولید تجدیدپذیر، باتری و تبادل توان با شبکه اصلی است. در چنین مسئله‌ای، دقت داده‌ای به‌تنهایی کافی نیست و خروجی مدل باید با توازن توان و محدودیت‌های تجهیزات سازگار باشد.

ایده اصلی شبکه فیزیک‌آگاه آن است که قوانین حاکم بر سامانه در تابع هزینه آموزش وارد شوند. این ایده از چارچوب عمومی رئیسی و همکاران برای حل و شناسایی معادلات دیفرانسیل الهام گرفته است \pcite{raissi2019pinns}. برای صورت‌بندی بخش ریزشبکه نیز از ساختار مدیریت انرژی، ذخیره‌ساز و تبادل توان در پژوهش اولادجو و فولی استفاده شد \pcite{oladejo2019microgrid}. بر این اساس، مؤلفه‌های توازن توان، پویایی وضعیت شارژ، حدود بهره‌برداری و هدف اقتصادی در یک تابع هزینه ترکیبی قرار گرفتند.

چارچوب ارائه‌شده یک طرح قابل پیاده‌سازی است و درباره برتری عددی مدل ادعایی ندارد. نتیجه‌گیری تجربی تنها پس از آماده‌سازی داده، اجرای کد و مقایسه منصفانه مدل فیزیک‌آگاه با مدل‌های پایه امکان‌پذیر است. این تفکیک میان طراحی روش و نتیجه تجربی برای حفظ اعتبار علمی پروژه ضروری است.

\section{پیشنهادها برای ادامه کار}
مسیرهای زیر می‌توانند در ادامه پروژه بررسی شوند:
\begin{enumerate}
  \item تکمیل مجموعه‌داده یا ساخت شبیه‌ساز معتبر برای ریزشبکه؛
  \item بررسی راهبردهای مختلف برای وزن‌دهی مؤلفه‌های تابع هزینه؛
  \item مقایسه شبکه فیزیک‌آگاه با شبکه عصبی صرفاً داده‌محور و یک روش بهینه‌سازی کلاسیک؛
  \item افزودن عدم قطعیت تولید خورشیدی، قیمت و بار؛
  \item استفاده از قید سخت به جای جریمه نرم برای متغیرهای حساس؛
  \item توسعه مدل برای حالت جزیره‌ای یا چند ریزشبکه متصل؛
  \item ترکیب نظریه بازی با مدل فیزیک‌آگاه برای تخصیص منصفانه سود؛
  \item بررسی مدل‌های زمانی مانند شبکه بازگشتی یا ترنسفورمر فیزیک‌آگاه؛
  \item ارزیابی زمان استنتاج برای کاربردهای برخط؛
  \item انتشار کد و تنظیمات آزمایش برای افزایش تکرارپذیری.
\end{enumerate}

\section{سخن پایانی}
شبکه‌های عصبی فیزیک‌آگاه پلی میان مدل‌سازی ریاضی و یادگیری از داده ایجاد می‌کنند. ارزش این رویکرد در مسئله ریزشبکه زمانی روشن می‌شود که هم معیارهای یادگیری و هم قوانین فیزیکی با دقت تعریف و ارزیابی شوند. نسخه نهایی پروژه باید این چارچوب را با داده و آزمایش واقعی تکمیل کند و محدودیت‌های مشاهده‌شده را نیز در کنار نقاط قوت گزارش دهد.
